% Checked necessary-condition supplement to Finding A44.
\paragraph{Supplement to Finding A44: odd auxiliary-prime towers.}

This note continues the controlled-tower test after Finding~A44.
The pro-two conclusion does not extend by the same argument to an
odd auxiliary-prime tower.  There is, however, a precise reduction
of the orders which that method would have to produce.  This is
an analysis of the remaining gate, not a proof that the required
orders occur.

Let $A/\mathbf F_q$ be an abelian surface, where $q$ is a power
of two, and write its Frobenius polynomial as
\[
 P_A(T)=T^4-aT^3+bT^2-qaT+q^2.
\]
Put
\[
 D_+=(b-q^2-1)(q+1)^2-qa^2\ne0.
 \tag{OT.1}
\]
The nonvanishing is the argument already proved in A44: equality
would give a factor $T^2-sT+1$ of $P_A$ over $\mathbf Q$, with
$s=a/(q+1)$, although the product of the absolute values of two
Weil roots is $q>1$.
Fix an odd prime $r$, and put
\[
 K_r=\bigcup_{m\ge0}\mathbf F_{q^{r^m}},\qquad
 S_0=\{\hbox{odd primes dividing }
           r\,\#A(\mathbf F_q)(q-1)D_+\}.
 \tag{OT.2}
\]
Then every odd prime $\ell\notin S_0$ occurring in the order of a
point of $A(K_r)$ satisfies $\ell\equiv1\pmod r$.
More precisely, if $m\ge1$ is the first index for which
$\ell\mid\#A(\mathbf F_{q^{r^m}})$, then
\[
 \ell\equiv1\pmod{r^m},
 \qquad P_A(T)\bmod\ell
 \text{ has a root in }\mathbf F_\ell
 \text{ of exact order }r^m.
 \tag{OT.3}
\]

\begin{proof}
Since $\ell\ne r,2$, the polynomial $T^{r^m}-1$ is separable
modulo $\ell$.  The first-appearance argument from A44 therefore
gives a Frobenius eigenvalue $\lambda$ on $A[\ell]$ of exact
order $r^m$.  Its degree over $\mathbf F_\ell$ is
\[
                 d=\operatorname{ord}_{r^m}(\ell)\le4.
\]
The multiset of roots of $P_A\bmod\ell$ is stable under the
involution $\alpha\mapsto q/\alpha$.  Its distinct Frobenius
orbits are either equal or disjoint, including when the
characteristic polynomial has repeated factors.

Degree three is impossible.  The orbit of three roots of the
irreducible factor containing $\lambda$ and its reciprocal orbit
cannot be disjoint inside a polynomial of degree four.  They must
be the same orbit.  The reciprocal involution commutes with its
three-cycle Frobenius permutation.  The centralizer of that
three-cycle on its three elements is cyclic of order three, so
the commuting involution is the identity.  This would give
$\lambda^2=q$, which has degree at most two over $\mathbf F_\ell$,
a contradiction.

If $d=2$, the cyclicity of
$(\mathbf Z/r^m\mathbf Z)^\times$ gives
$\ell\equiv-1\pmod{r^m}$, so
$\lambda^{\ell+1}=1$.  Its irreducible quadratic is therefore
$g=T^2-sT+1$ for some $s\in\mathbf F_\ell$.
If its reciprocal orbit coincides with its orbit, the reciprocal
involution cannot be the identity: that would put
$\lambda^2=q$ in $\mathbf F_\ell$, and since $\lambda$ has odd
order, $\lambda$ itself would belong to $\mathbf F_\ell$.
It must be the nontrivial Frobenius permutation, and consequently
$q=\lambda^{\ell+1}=1$ modulo $\ell$.
Otherwise the two quadratic orbits are disjoint and exhaust the
four roots, giving
\[
 P_A(T)=(T^2-sT+1)(T^2-qsT+q^2)\pmod\ell.
\]
Comparing coefficients without division gives $\ell\mid D_+$.
Both possibilities are excluded by $\ell\notin S_0$.

If $d=4$, its orbit exhausts all four roots.  The reciprocal
involution commutes with the four-cycle Frobenius permutation.
It is not the identity, again because that would give
$\lambda^2=q$ and degree at most two.  It is therefore the square
of that permutation.  Since an element of order four in the
cyclic group $(\mathbf Z/r^m\mathbf Z)^\times$ squares to $-1$,
\[
 q=\lambda^{\ell^2+1}=1\pmod\ell,
\]
which is also excluded.  Only $d=1$ remains, proving (OT.3).
For $r=3$ the degree-four case cannot arise; the same proof
otherwise applies without a restriction $r\ge5$.
\end{proof}

The whole-order condition is substantially stronger than
eligibility of its individual prime divisors.  Call an odd
integer $n>1$ compatible when
\[
                   \langle2,-1\rangle
                      =(\mathbf Z/n\mathbf Z)^\times.
 \tag{OT.4}
\]
If a point of $A(K_r)$ has compatible order $n$, at most one
prime divisor of $n$ lies outside $S_0$.  If such a prime occurs,
then $v_r(n)\le1$.  In particular, the sparse cofactor supplied by
Poonen's tower proof cannot retain two new primary components
and still satisfy (OT.4).

\begin{proof}
The Sylow-$r$ subgroup of a group generated by $2$ and $-1$ is
cyclic: the subgroup generated by $2$ has index at most two,
and $r$ is odd.  Each distinct prime $\ell\notin S_0$ dividing
$n$ supplies a nontrivial Sylow-$r$ factor in
$(\mathbf Z/\ell^{v_\ell(n)}\mathbf Z)^\times$, since
$r\mid\ell-1$.  Two such factors would give a noncyclic product.
If $r^2\mid n$, the factor for $r^{v_r(n)}$ supplies another
nontrivial Sylow-$r$ factor, proving the second assertion.
The same argument also excludes any other cofactor prime $s$
with $r\mid s-1$ when a prime outside $S_0$ already occurs.
\end{proof}

There is an explicit finite list of possible cofactors.  Let
\[
 e=\max_{s\in S_0\setminus\{r\}}
       v_r\!\left(\prod_{j=1}^4(s^j-1)\right),
 \qquad
 N_0=\#\bigl(A(\mathbf F_{q^{r^e}})\{S_0\setminus\{r\}\}\bigr),
 \tag{OT.5}
\]
where an empty maximum is zero, and the final notation denotes
the order of the indicated primary subgroup.  Then
\[
 A(K_r)\{S_0\setminus\{r\}\}
  =A(\mathbf F_{q^{r^e}})\{S_0\setminus\{r\}\}.
 \tag{OT.6}
\]
Indeed, a point of $s$-power order, $s\ne r$, has a Frobenius
orbit of $r$-power length dividing the order of
$\mathrm{GL}_4(\mathbf Z/s^j\mathbf Z)$ for some $j$.
Its $r$-part is bounded by the product in (OT.5), independently
of $j$, because the reduction kernel is an $s$-group.  This is
Poonen's complete Lemma~6.1 argument and the finite-subgroup
conclusion of Lemma~6.2, with the field bound made explicit
\cite[Lemmas~6.1--6.2]{Poonen2005}.

Also, (OT.4) forces $n$ to have at most two distinct odd prime
factors: its unit group has $2$-rank equal to their number,
whereas two generators give rank at most two.  Thus any compatible
$n$ containing a prime outside $S_0$ has the form
\[
 n=c\ell^a,\quad a\ge1,\quad \ell\notin S_0,\qquad
 c\in\mathcal M,
 \tag{OT.7}
\]
where the following possibly redundant set is finite and
computable:
\[
 \mathcal M=\{1,r\}\ \cup
 \{s^j:s\in S_0\setminus\{r\},\ r\nmid s-1,
                         \ 1\le j\le v_s(N_0)\}.
 \tag{OT.8}
\]
One must still test (OT.4) at the actual integer in (OT.7).
The inequalities and the list are necessary conditions, not
sufficient conditions for a theta incidence.

For a genus-two curve in characteristic two, take
$A=J_{C^{(1)}}$ and the actual Raynaud support
$\Theta_C=C^{(1)}-B_C$.  Finding~A19 makes its subset of points
supported on the fixed odd set $S_0$ finite, with an effective
finite-field enumeration bound.  Consequently an unbounded
compatible theta supply from this odd tower would require
infinitely many incidences of the form (OT.7), with $c$ in the
fixed finite list.  The pure-prime-power case $c=1$ remains among
them.  This conclusion does not say that no such incidences
exist or that an arbitrary individual curve lacks a witness.

\paragraph{Exact location of the unresolved inference.}
In Poonen's proof of Lemma~6.4, a large finite-prime initial set
is removed and the remaining subgroup $H$ has uniformly bounded
index; a translated Lang pullback then produces a point in
$a+H$.  This controls the total size of the removed primary
part, but allows the point's residual component to use several
new primes.  In the tower above, two such primes already
contradict (OT.4).  Replacing $H$ by the union of its individual
primary subgroups would impose the needed restriction, but
that union need not be a subgroup and does not define the same
Lang quotient.  The bounded-index argument therefore stops at
this step.  The density-zero conclusion of Poonen's
Lemma~6.3 does not remove this issue
\cite[Lemmas~6.3--6.6]{Poonen2005}.

Nor is the surviving congruence itself an arithmetic obstruction
to eligibility: for $r=3$, the primes $7$ and $13$ are both
$1\bmod3$ and satisfy $\langle2,-1\rangle=U_\ell$.
Their product $91$ does not satisfy that equality, since its
unit group has a Sylow-three subgroup $C_3\times C_3$.
These are assertions about finite unit groups only; they do not
assert Frobenius eigenvalues or theta points on a chosen curve.

This probe sharpens the residual selection gate for the
controlled-tower construction.  It supplies neither a
prime-incidence theorem nor a resolution of a previously open
universal prescribed-point row.  No density statement about
Artin primes alone is used to infer incidence.
